\chapter{物理専攻者の学習導線}
数学・物理の既知知識を、工学と地球科学の観測・設計・判断へ接続する。
新規用語の数、推論段階、計算負荷、図表の読解を区別して、自分にとって必要な準備を見積もる。
\section{公式問題群と学問体系}
基礎の公式5群は出題上の区分として保持する。一方、勉強するときは線形代数、振動、材料力学、数値解析などの学習分野へ分けて考える。
「固有」という共通語があっても、行列の性質と構造物の固有振動は同じ学習項目にはまとめない。
\section{初めて学ぶ分野の読み方}
情報分野では、状態と操作を明示することが物理の状態変数と時間発展の考え方に対応する。
化学分野では、反応式だけでなく溶液条件と、操作によって残る成分を追う。
地球科学では、一般理論と、実際の自然試料・観測データからの解釈を分ける。
制度・倫理では、出題当時の記述、公式正答、現行資料の説明をそれぞれ記録する。
\section{問題を解いた後}
問題幹の否定、単位、図の境界条件、選択肢番号の対応まで戻って確認する。
同じ解法の問題でも数値や選択肢順が変われば、正答番号を記憶する方法は通用しない。
未確認の類似候補を再出題と確定せず、差分を比較して承認した関係だけを復習に使う。
